\documentclass[12pt,a4paper]{report}

\usepackage[utf8]{inputenc}
\usepackage{graphicx}
\usepackage{geometry}
\geometry{a4paper, margin=1.2in}
\usepackage{setspace}
\usepackage{titlesec}
\usepackage{hyperref}
\usepackage{multirow}
\usepackage{array}
\usepackage{xcolor}

% Formatting chapters to look like the ones in the PDF (numbered linearly or properly scaled)
\titleformat{\chapter}[display]
  {\normalfont\bfseries}{}{0pt}{\Large}
\titlespacing*{\chapter}{0pt}{-20pt}{20pt}

\begin{document}
\setstretch{1.2} % Line spacing

% ---------------------------------------------------------
% TITLE PAGE
% ---------------------------------------------------------
\begin{titlepage}
    \begin{center}
        \vspace*{1cm}
        
        \textbf{\Large "Optimizing Cloud-Native Workflow Orchestration: A Caching-Centric Approach for DAG Execution"}
        
        \vspace{1.5cm}
        
        \textit{B.Tech. Project (CS400) Report}\\
        \textit{submitted in partial fulfilment of the requirements for the award of the degree of}\\
        
        \vspace{0.8cm}
        \textbf{\large Bachelor of Technology}
        
        \vspace{1.5cm}
        
        \textit{by}\\
        \vspace{0.8cm}
        \textbf{\Large Pratik Suman Surin}\\
        \vspace{0.3cm}
        (Roll No: 2201149)
        
        \vspace{2.5cm}
        
        \textit{Under the supervision of}\\
        \textbf{\Large Dr. Shubha Brata Nath}
        
        \vfill
        
        % Uncomment and add your logo image here
        % \includegraphics[width=0.3\textwidth]{logo.png} \\
        \vspace{1cm}
        
        \textbf{\large Department of Computer Science and Engineering}\\
        \textbf{\large Indian Institute of Information Technology Guwahati}\\
        \textbf{Guwahati, India}\\
        \textbf{November 2025}
        
    \end{center}
\end{titlepage}

% ---------------------------------------------------------
% DECLARATION
% ---------------------------------------------------------
\newpage
\thispagestyle{empty}
\begin{center}
    \textbf{\Large DECLARATION}
\end{center}
\vspace{1cm}
I hereby \textit{declare} this thesis entitled \textbf{"Optimizing Cloud-Native Workflow Orchestration: A Caching-Centric Approach for DAG Execution"} that is being submitted to \textbf{Indian Institute of Information Technology Guwahati, Assam}, in partial fulfilment for the requirements of the award of degree of \textbf{Bachelor of Technology} in \textbf{Computer Science and Engineering} in the department of \textbf{Computer Science and Engineering}, is a \textit{genuine report of the work carried out by me}. The material contained in this report has not been submitted at any other University or Institution for the award of any degree.

\vspace{3cm}

\begin{flushright}
    \textbf{Pratik Suman Surin} \\
    \textbf{Roll No: 2201149} \\
    \vspace{1cm}
    ......................................................... \\
\end{flushright}

\begin{flushleft}
    Place: IIIT Guwahati, Assam \\
    Date: December 2, 2025 \\
\end{flushleft}
\vspace{-1.5cm}
\begin{flushright}
    (Signature of the Student) \\
    Computer Science and Engineering
\end{flushright}


% ---------------------------------------------------------
% ABSTRACT
% ---------------------------------------------------------
\newpage
\pagenumbering{roman}
\addcontentsline{toc}{chapter}{Abstract}
\begin{center}
    \textbf{\Large Abstract}
\end{center}
\vspace{1cm}
Distributed workflow orchestration forms the backbone of many modern data engineering systems, yet traditional approaches often encounter substantial performance bottlenecks due to network-dependent artifact transfers. These limitations become especially pronounced in data-intensive pipelines, where repeated communication with remote storage introduces unnecessary delays and reduces the overall efficiency of workflow execution.

This project presents, a Kubernetes-native Directed Acyclic Graph (DAG) orchestration simulation designed to address these challenges by minimizing execution overhead through improved data locality. Helios introduces a node-level caching architecture that leverages local RAM on the node to bypass remote storage for intermediate data movement, significantly streamlining artifact access within the pipeline.

In addition, the platform structures task execution in a way that naturally promotes co-location of dependent operations on the same Kubernetes nodes. This design maximizes the use of node-local caches, minimizes remote I/O, and improves overall system responsiveness by ensuring that data movement patterns are inherently aligned with where tasks are executed.

Experimental evaluations indicate that this integrated caching and scheduling strategy results in noticeably faster artifact retrieval and a marked reduction in total pipeline latency during warm executions. The outcomes highlight the effectiveness of combining node-local caching with advanced Kubernetes scheduling logic to enhance throughput, reduce network congestion, and optimize performance in cloud-native workflow engines.

% ---------------------------------------------------------
% TABLE OF CONTENTS
% ---------------------------------------------------------
\newpage
\tableofcontents

% ---------------------------------------------------------
% MAIN DOCUMENT CHAPTERS
% ---------------------------------------------------------
\newpage
\pagenumbering{arabic}

\chapter{1 \hspace{0.2cm} Introduction}
Distributed systems and cloud-native architectures have become essential in modern data processing pipelines...
% Add introduction paragraphs here

\section*{1.1 Overview}
\addcontentsline{toc}{section}{1.1 Overview}
Cloud orchestration is the technology used to manage complex workflows, commonly represented as \textit{Directed Acyclic Graphs} (DAGs). In such systems, a central controller ensures that tasks execute in the correct order while efficiently utilizing resources across a cluster of machines.

In typical cloud environments, the nodes responsible for computation are separate from the systems responsible for storage...

\section*{1.2 Challenges}
\addcontentsline{toc}{section}{1.2 Challenges}
\subsection*{1.2.1 Network Overhead and Latency}
\addcontentsline{toc}{subsection}{1.2.1 Network Overhead and Latency}
There is a significant challenge in cloud orchestration arises from the inherent limitations of network communication...

\subsection*{1.2.2 Inefficient Resource Utilization}
\addcontentsline{toc}{subsection}{1.2.2 Inefficient Resource Utilization}
Slow or inconsistent data transfer results in substantial idle time for compute nodes...

\subsection*{1.2.3 Unpredictable Execution Times}
\addcontentsline{toc}{subsection}{1.2.3 Unpredictable Execution Times}
Network variability significantly impacts workflow predictability...

\subsection*{1.2.4 Difficulty in Error Tracking}
\addcontentsline{toc}{subsection}{1.2.4 Difficulty in Error Tracking}
Error diagnosis in distributed cloud workflows presents significant difficulties...

\section*{1.3 Motivation}
\addcontentsline{toc}{section}{1.3 Motivation}
The primary motivation for this research is due to the critical inefficiencies present in modern cloud infrastructure...

\section*{1.4 Problem Statement}
\addcontentsline{toc}{section}{1.4 Problem Statement}
The central problem addressed by this research is the significant performance penalty caused by the physical separation of computation and storage in distributed cloud environments...

\vspace{0.5cm}
\noindent \textbf{1. The ``Data Movement'' Bottleneck}\\
\begin{itemize}
    \item \textbf{Cold Start Data Problem:} Every task begins with an empty workspace...
    \item \textbf{Network Congestion:} When multiple tasks simultaneously request large files...
\end{itemize}

\vspace{0.5cm}
\noindent \textbf{2. Inefficiency of ``Blind'' Scheduling}\\
\begin{itemize}
    \item \textbf{Random Task Assignment:} Traditional schedulers don't care of where the data is stored...
    \item \textbf{Redundant Data Transfers:} A scheduler may assign a task to Node A...
\end{itemize}

% ---------------------------------------------------------
\chapter{2 \hspace{0.2cm} Related Works}
\begin{itemize}
    \item \textbf{High-Throughput Orchestration: ArmoniK}\\
    The ArmoniK framework is an advanced, open-source orchestrator designed for High-Throughput Computing...
    \item \textbf{Data-Locality Aware Workflow Scheduling (D-LAWS)}\\
    D-LAWS addresses the challenge of minimizing data movement in scientific workflows...
    \item \textbf{Dependency-Aware Caching in DAG Systems}\\
    Research in frameworks like Apache Spark shows that traditional caching policies...
\end{itemize}

% ---------------------------------------------------------
\chapter{3 \hspace{0.2cm} System Architecture}
The system employs a cloud-native, microservices architecture designed to enforce data locality and minimize network latency...

% Uncomment to place images
% \begin{figure}[h]
% \centering
% \includegraphics[width=0.9\textwidth]{architecture.png}
% \caption{Helios Workflow Orchestration System Architecture}
% \end{figure}

\section*{3.1 Control Plane: The Orchestration Framework}
\addcontentsline{toc}{section}{3.1 Control Plane: The Orchestration Framework}
\begin{itemize}
    \item \textbf{Scheduler Service (Java/Spring Boot):}
    \begin{itemize}
        \item Parses user-submitted DAGs.
        \item Tracks task dependencies.
        \item Manages workflow lifecycle.
    \end{itemize}
    \item \textbf{Kubernetes API Server:}
    \begin{itemize}
        \item Interfaces between the Scheduler and cluster infrastructure.
        \item Embeds Node Affinity rules for data-aware placement.
    \end{itemize}
    \item \textbf{RabbitMQ (Message Broker):}
\end{itemize}

\section*{3.2 Data Plane: Persistence and Storage}
\addcontentsline{toc}{section}{3.2 Data Plane: Persistence and Storage}
\begin{itemize}
    \item \textbf{Redis (DAG State Store):}
    \item \textbf{MinIO (Artifact Storage):}
\end{itemize}

\section*{3.3 Execution Plane: Data-Aware Execution}
\addcontentsline{toc}{section}{3.3 Execution Plane: Data-Aware Execution}
\begin{itemize}
    \item \textbf{Task Pod (K8s Job):}
    \item \textbf{Node Cache (HostPath/RAM):}
    \item \textbf{Init Container (Smart Downloader):}
    \item \textbf{Main Container (Smart Uploader):}
\end{itemize}

% ---------------------------------------------------------
\chapter{4 \hspace{0.2cm} Algorithms}
The core of this system relies on two algorithms coupled through the central Redis state store:

\section*{4.1 Algorithm 1: Data Gravity Scheduling Heuristic}
\addcontentsline{toc}{section}{4.1 Algorithm 1: Data Gravity Scheduling Heuristic}

The scheduling decision is formalized as a resource placement problem...
\vspace{0.3cm}

\noindent \textbf{Input:} Target Task ($\mathcal{T}$), Workflow ID ($\mathcal{D}$) \\
\textbf{Output:} Preferred Execution Node ($\mathcal{N}^*$) or Null

\begin{table}[h]
\centering
\begin{tabular}{|c|c|p{9cm}|}
\hline
\textbf{Step} & \textbf{Operation} & \textbf{Description} \\
\hline
1 & Initialization & Initialize map to track data size scores per node ($Scores \leftarrow \{\}$) \\
\hline
2 & Dependency Check & Skip optimization if task has no dependencies. \\
\hline
3 & Calculate Gravity & For each predecessor task... \\
\hline
\end{tabular}
\end{table}

\section*{4.2 Algorithm 2: Node-Level Caching Logic}
\addcontentsline{toc}{section}{4.2 Algorithm 2: Node-Level Caching Logic}
The caching mechanism is a concurrent, local persistence logic implemented in the Init Container...

% ---------------------------------------------------------
\chapter{5 \hspace{0.2cm} Experimentation}
\section*{5.1 Experimental Setup}
\addcontentsline{toc}{section}{5.1 Experimental Setup}
\begin{itemize}
    \item \textbf{Simulated Environment:} Execution utilized the microservices System over a network...
    \item \textbf{Workloads:} Workflows were constructed using a classic I/O-bound pattern...
\end{itemize}

\section*{5.2 Results}
\addcontentsline{toc}{section}{5.2 Results}
% Example Table
\begin{table}[h]
\centering
\begin{tabular}{|c|c|c|c|c|}
\hline
\textbf{Artifact Size} & \textbf{Non-Cached (s)} & \textbf{Cached (s)} & \textbf{Time Saved (s)} & \textbf{Reduction (\%)} \\
\hline
200 MB & 60.00 & 41.12 & 18.88 & 31.47\% \\
\hline
100 MB & 53.17 & 37.90 & 15.27 & 28.72\% \\
\hline
\end{tabular}
\caption{Comparison of Workflow Durations Under Cached and Non-Cached Conditions}
\end{table}

\section*{5.3 Visualizations}
\addcontentsline{toc}{section}{5.3 Visualizations}
\section*{5.4 Key Findings}
\addcontentsline{toc}{section}{5.4 Key Findings}

% ---------------------------------------------------------
\chapter{6 \hspace{0.2cm} Conclusion}
This research introduced a framework for cloud orchestration designed to overcome the critical performance bottleneck associated with data movement in disaggregated computing environments...

% ---------------------------------------------------------
\renewcommand{\bibname}{References}
\begin{thebibliography}{9}
\addcontentsline{toc}{chapter}{References}
\bibitem{choi} 
Choi, J., Adufu, T., \& Kim, Y. (2017). Data-Locality Aware Scientific Workflow Scheduling Methods in HPC Cloud Environments. International Journal of Parallel Programming, 45(5), 1128–1141.

\bibitem{zhao}
Zhao, Y., Dong, J., Liu, H., Wu, J., \& Liu, Y. (2021). Performance Improvement of DAG-Aware Task Scheduling Algorithms with Efficient Cache Management in Spark. Electronics, 10(16), 1874.

\bibitem{delamea}
Q. Delamea et al., "Cloud-Agnostic Serverless Platform for Fault-Tolerant Execution of Dynamic Task Graphs," 2025 IEEE Cloud Summit.
\end{thebibliography}

\end{document}
